\newpage

\section{\huge \textcolor{blue}{二阶线性微分方程}}
\subsection{\LARGE 二阶常系数线性微分方程}
对与二阶常系数线性微分方程:
\[ y'' + A y' + B y = 0 \]
先给出通解:
\[ y = c_1 y_1 + c_2 y_2 \]

$ y_1, y_2 $ 为通解的2个解, 
即可将一次求解任务拆解为两个单次求解任务
\subsubsection{\large 阻尼运动}
下以简单阻尼运动模型为例:
\begin{figure}[htbp]
    \centering
\includegraphics[width=0.6\textwidth]{C:/Users/37/Desktop/Notes on Common Calculus Equations/Notes and pictures/c8.1.1.png}
    \caption{阻尼模型}
    \label{fig:c8.1.1}
\end{figure}
\[ m \ddot{x} + c \dot{x} + k x = 0 \]
\[  \ddot{x} + \frac{c}{m} \dot{x} + \frac{k}{m} x = 0 \]
下即针对解此类方程进行研究:\\
由通解形式可知, 得出的两解应线性无关\\
先假设其中一解为:
\[ y =  e^{\gamma t}  \]
有:
\[  y' = \gamma e^{\gamma t} \]
\[  y'' = \gamma^2 e^{\gamma t} \] 
\[\therefore \hspace*{2em}  \gamma^2 e^{\gamma t} + A \gamma e^{\gamma t} + B e^{\gamma t} = 0 \]
\[\therefore \hspace*{2em}  \gamma^2 + A \gamma + B = 0 \]

\subsubsection{\large ODE特征方程}
{\Large \[ \gamma^2 + A \gamma + B = 0 \] }
由上得此方程, 即为二阶方程的ODE特征方程\\
解得:
\[  \gamma = - \frac{A}{2} \pm \sqrt{\frac{A^2}{4} - B}\]
为形式更加和谐, 将方程系数改写为:
\[  \gamma^2 + 2 p \gamma + \omega ^2 = 0 \] 
解得:
\[  \gamma = - p \pm \sqrt{p^2 - \omega ^2} \]
由解形式, 可知根有3种情况:


case 1: $p^2 - \omega ^2 > 0$\\
根为实数且不相等\\
解为:
\[ y = c_1 e^{r_1 t}  + c_2 e^{r_2 t}  \]
若给出初值, 如$y(0) = 1, y'(0) = 0$, 可求得$c_1, c_2 $值， 得最后结果. 


case 2: $p^2 - \omega ^2 = 0$\\
根为重根\\
方程为:
\[ \gamma^2 + 2 p \gamma + p ^2 = 0 \]
\[ \gamma = - p \]
现得一解$e^{-pt}$, 无另根来对应另一解, 故现在问题为找另根\\
有下:
\[ y'' + p y' + p^2 y = 0 \]
具有如下形式解:
\[ y = y_1 u \]
以 $y_1$为因子, $y = y_1 u$是方程的另一个线性无关解(u非常数)\\
解:
\begin{align*}
    y  &= e^{-pt} u\\
    y' &= -p e^{-pt} u + e^{-pt} u' \\
    y''&= p^2 e^{-pt} u -2p e^{-pt} u' + e^{-pt} u'' 
\end{align*}

将上式分别乘$p^2, p, 1$后相加, 得:
\[ e^{-pt} u'' = 0 \]
\[ \therefore \hspace*{1em} u' = 0 \]
\[ u = c_1 t+ c_2  \]
即找一解即可, 可取 $u = t$
\[ \therefore \hspace*{1em}  y = e^{-pt} t \]


case 3: $p^2 - \omega ^2 < 0$\\
根为复数\\
解得:
\[ \gamma = a \pm b i \]
\[ \therefore y = e^{ a \pm b i} \]
选取 $ a + b i $, 即可:
\begin{align*}
    y  &= e^{at + b t i} \\
       &= e^{at} (c_1 \cos bt + c_2 \sin bt)\\
       &= c_1 y_1 + c_2 y_2 
\end{align*}

即得完整解, 那为什么不用 $a - b i$呢?\\
观察知道, 因
\[ y_1 = e^{at} \cos bt = e^{at} \cos (-bt) = y_1 \]
\[ y_2 = e^{at} \sin bt = - e^{at} \sin (-bt) = -y_2 \]
只有其一系数改变符号, 其余不变, 则无需多余考虑\\
那能否直接用如下表示解?
\[ y = C_1 e^{(a + b i) t} + C_2 e^{(a - b i) t} \]
其中$C_1, C_2$为复系数\\
求证2种方法答案一致有2种方法, 如下:
\begin{enumerate}
    \item 完全展开， 乘法分配, 取实部.(不展开说明)
    \item 利用 $u + i v$, 可令$i$为$-i$法找到实部
\end{enumerate}

具体操作如下:
\[ C_1 e^{(a + b i) t} + C_2 e^{(a - b i) t} \]
$i \rightarrow -i $
\[ \bar{C_1} e^{(a - b i) t} + \bar{C_2} e^{(a + b i) t} \]
$\Rightarrow $
\[ C_1 = \bar{C_2} , C_2 = \bar{C_1} \]
则令
\[ y = (c + di) e^{(a + b i) t} + (c - di) e^{(a - b i) t} \]
变形:
\[ y = e^{at} [c (e^{ibt} + e^{-ibt}) + i d (e^{ibt} - e^{-ibt})] \]
利用如下欧拉公式:
\[ \cos x = \frac{1}{2} (e^{ix} + e^{-ix}) \]
\[ \sin x = \frac{1}{2i} (e^{ix} - e^{-ix}) \]
得:
\[ y = e^{at} (2c \cos bt + 2d \sin bt) \]
即与前只选取$a + bi$来求解的形式一致

\subsection{\LARGE 二阶齐次线性微分方程}
\subsubsection{\large 解理论}
{\Large \[ y'' + p(x) y' + q(x) y = 0 \] }
由前文讨论, 知求解方式为:\\
找到线性无关的$y_1, y_2$两解(不考虑$0$解), 进行线性组合
即为完整解, 自然引出下2个问题:


Q1: $c_1 y_1 + c_2 y_2$为什么是解？\\
下引入叠加原理:\\
若$y_1, y_2$是线性齐次ODE的解\\
$\Longrightarrow $  $c_1 y_1 + c_2 y_2$也是解\\
为说明此原理, 介绍线性算子:\\
对原方程:
\[ y'' + p y' + q y = 0 \]
变形:
\[ D^2 y + p D y + q y = 0 \]
\[ (D^2 + p D + q) y = 0 \]
记$L = D^2 + p D + q$ 为线性算子
\[ L y = 0 \]
对于此将线性运算打包成一个算子, 称为线性算子\\
此处线性算子为: $L = D^2 + p D + q$\\
下讨论$L$为一般的线性算子时, 存在性质:
\begin{enumerate}
    \item $L(u_1 + u_2) = L(u_1) + L(u_2)$
    \item $L (cu) = c L(u)$
\end{enumerate}
由此即可得Q1的解:\\
若$y_1,y_2$为$L$的解, 则
\begin{align*}
    L(c_1 y_1 + c_2 y_2) &= L(c_1 y_1) + L(c_2 y_2) \\
         &= c_1 L(y_1) + c_2 L(y_2)\\
         &= 0 + 0 \\
         &= 0
\end{align*}

即说明$c_1 y_1 + c_2 y_2$是$L$的解.\\[2em]
\hspace*{2em} Q2: 为什么$c_1 y_1 + c_2 y_2$是所有解?\\
为解决Q2, 从初值的角度出发,已知解与任意初值
是否能找到一对$c_1, c_2$满足, 又是否唯一?\\
下仅做探讨:\\
现有:
\[
\left\{
\begin{aligned}
    y(x_0) &= a \\
    y'(x_0) &= b
\end{aligned}
\right.
\]
即:
\[
\left\{
\begin{aligned}
    y = c_1 y_1 + c_2 y_2 \\
    y' = c_1 y'_1 + c_2 y'_2
\end{aligned}
\right.
\]
\[
\left\{
\begin{aligned}
    c_1 y_1 + c_2 y_2 = a \\
    c_1 y'_1 + c_2 y'_2 = b
\end{aligned}
\right.
\]
由代数知识:\\
若$c_1 , c_2$可解, 满足:
\[
\begin{vmatrix}
    y_1 & y_2 \\
    y'_1 & y'_2
\end{vmatrix}_{x_0}
\neq 0
\]
补充:\\
将形如上式的行列式称为\\
Wronskian(朗斯基行列式):
\[
W(y_1,y_2;x) = 
\begin{vmatrix}
    y_1 & y_2 \\
    y'_1 & y'_2
\end{vmatrix}
\]
且有:\\
\[ W(y_1,y_2;x) \equiv 0 \hspace*{1em} or \neq 0 \text{(for all x)} \]
补充到此, 继续讨论:\\
当写出\{$c_1y_1 + c_2y_2$\}时, 即是宣称所以解\\
假设\{$c_1'u_1 + c_2'u_2$\}是另一对线性无关组:
\[
\left\{
\begin{aligned}
    u_1 = \hat{c_1} y_1 + \hat{c_2} y_2  \\
    u_2 = \hat{c_1} y_1 + \hat{c_2} y_2
\end{aligned}
\right.
\]
为有更佳的初值， 自然正交化:
\[
\left\{
\begin{aligned}
    Y_1(0) = 1 ,  Y_1'(0) = 0 \\
    Y_2(0) = 0 ,  Y_2'(0) = 1
\end{aligned}
\right.
\]
下给出几个例子:\\
eg1:
\[ y'' + y = 0 \]
由初值解下:
\[
\left\{
\begin{aligned}
    y_1 = \cos t = Y_1  \\
    y_2 = \sin t = Y_2
\end{aligned}
\right.
\]
eg2:
\[ y'' - y = 0 \]
由初值解下:
\[
\left\{
\begin{aligned}
    y_1 = c_1 e^x + c_2 e^{-x} \\
    y_2 = c_1 e^x - c_2 e^{-x}
\end{aligned}
\right.
\]
得:
\[
\left\{
\begin{aligned}
    Y_1 = \frac{(e^x +  e^{-x})}{2} = \cosh x  \\
    Y_2 = \frac{(e^x -  e^{-x})}{2} = \sinh x
\end{aligned}
\right.
\]
观察, 有如下定理:\\
若$Y_1, Y_2$在$0$处正交化, 初值问题为:
\[
\left\{
\begin{aligned}
    y(0) = y_0  \\
    y'(0) = y_0'
\end{aligned}
\right.
\]
解为
\[ y = y_0 Y_1 + y_0' Y_2 \]
\subsubsection{\large 存在和唯一性定理}
由上简单的讨论, 引入存在和唯一性定理:\\
若有方程如下:
\[ y'' + p(x)y' + q(x)y = f(x) \]
且$p, q$为连续函数, $y$为可导函数,
那么有且仅有唯一解满足给定初值条件:
\[
\left\{
\begin{aligned}
    y(0) = y_0  \\
    y'(0) = y_0'
\end{aligned}
\right.
\]
\[ y = y_0 y_1 + y_0' y_2 \]

\subsection{\LARGE 二阶非齐次线性微分方程}
\subsubsection{\large 解结构}
现考虑二阶非齐次线性微分方程:
{\Large \[ y'' + p(x)y' + q(x)y = f(x) \] }
对应的齐次方程(相伴齐次方程)为:
\[ y'' + p(x)y' + q(x)y = 0 \]
其解为:
\[ y_c = c_1 y_1 + c_2 y_2 \]
考虑非齐次方程的解结构, 方程记为:
\[ L y = f(x) \]
其中$L$为二阶线性算子\\
解形式:
\[ y = y_p + y_c \]
\[ y = y_p + c_1 y_1 + c_2 y_2 \]
其中$y_p$为$L y = f(x)$的特解, 
$y_c$为上述相伴齐次方程的解\\
为进一步讨论, 下提出2个命题并证明:\\
命题1: $y_p + c_1 y_1 + c_2 y_2$是解\\
证明:\\
代入方程:
\begin{align*}
L(y_p + c_1 y_1 + c_2 y_2) &= L(y_p) + L(c_1 y_1) + L(c_2 y_2) \\
    &= f(x) + 0 \\
L(y)    &= f(x)   
\end{align*}
证毕.\\[2em]
命题2: $y_p + c_1 y_1 + c_2 y_2$是唯一解\\
证明:\\
假设$u(x)$为解:
\begin{align*}
L(u) &= f(x) \\
L(y_p) &= f(x)\\
L(u - y_p) &= 0
\end{align*}
得, $u - y_p$是相伴齐方程的解, 
令$u - y_p = c_1 y_1 + c_2 y_2$\\
由存在和唯一性定理, 矛盾, 即无其它解.\\
证毕.\\
为了进一步讨论, 下回到一阶线性微分方程且只考虑
常系数的情况:
\[ y' + k y = q(t) \]
解:
\[ y = e^{-kt} \int q(t) e^{kt} dt + ce^{-kt} \]
观察, 套用前文结论:
\[
\left\{
\begin{aligned}
    &y_p = e^{-kt} \int q(t) e^{kt} dt\\
    &y_c = ce^{-kt} 
\end{aligned}
\right.
\]
且若$k > 0$, 还存在以下概念:
\[
\left\{
\begin{aligned}
    y_p = \text{稳态解}  \\
    y_c = \text{暂态解}
\end{aligned}
\right.
\]
思考:
\[ y'' + A y' + By = f(t) \]
当$t \to \infty$时, 在什么情况下, 暂态解为0$(\forall c_1, c_2)$?\\
若找到此情况, 则称这个ODE为稳定的.\\
从根开始考虑:
\begin{center}
    \begin{tabular}{c|c|c}
        \text{特征方程的根} & \text{解}                             & \text{趋于0的条件}  \\ \hline
        $r_1 \neq r_2$     & $c_1 e^{r_1 t} + c_2 e^{r_2 t} $      & $r_1,r_2 \hspace*{0.5em} \textcolor{red}{ < 0}$   \\ \hline
        $r_1 = r_2$        & $(c_1 + c_2 t) e^{r_1 t}$             & $r_1 = r_2 \hspace*{0.5em} \textcolor{red}{ < 0}$  \\ \hline
        $r = a \pm b i$    & $e^{a t} (c_1 \cos bt + c_2 \sin bt)$ & $a \hspace*{0.5em} \textcolor{red}{ < 0}$         \\
    \end{tabular}
\end{center}
观察得:\\
ODE稳定的条件是:
\begin{center}
所有特征根都具有负实部
\end{center}

\newpage
\subsection{\LARGE 非齐次 ODE}
\subsubsection{\large 指数型非齐次项}
继续观察下式方程:
\[ y'' + A y' + By = f(x) \]
若输入$f(x)$如下:
\[
\begin{aligned}
    f(x) &= e^{ax} \hspace*{2em} (a < 0 , often)\\
         &=  \left\{
                \begin{aligned}
                \sin \omega x  \\
                \cos \omega x
                \end{aligned}
                \right.\\
         &= \left\{
                \begin{aligned}
                &e^{ax} \sin \omega x  \\
                &e^{ax} \cosh \omega x
                \end{aligned}
                \right.   
\end{aligned}
\]
若 $\alpha \in \mathbb{C} $ \\
则, 以上都是$e^{\alpha X} = e^{(a + i \omega)x}$的特例\\
现将原式提取多项式算子:
\[ (D^2 + A D + B) y = f(x) \]
\[ P(D) y = f(x) \]
其中$p(D) = D^2 + A D + B$\\
同时存在如下结论:
\[ p(D) e^{\alpha x} = p(\alpha) e^{\alpha x} \]
代入, 即证.
\subsubsection{\large 指数输入定理}
关于方程:
\[ p(D) y = e^{\alpha x} \]
有特解:
\[ y_p = \frac{e^{\alpha x}}{p(\alpha)}  \hspace*{2em} p(\alpha) \neq 0 \]
代入, 即证.\\
思考, $p(\alpha) = 0$情况:\\
为解决此问题, 引入如下法则\\
指数移位法则:
\[ p(D) e^{\alpha x} u(x) = e^{\alpha x} p(D + \alpha) u(x)    \]
证明:\\
从最简单情况开始考虑:\\
$p(D) = D$
\begin{align*}
D(e^{\alpha x} u) &= e^{\alpha x} D u + a e^{\alpha x} u  \\
                &= e^{\alpha x} (D u + a u)  \\
                &= e^{\alpha x} (D + a) u 
\end{align*}
$p(D) = D^2$
\begin{align*}
D^2(e^{\alpha x} u) &= D(D(e^{\alpha x} u)) \\
                &= D(e^{\alpha x} (D + a) u)  \\
                &= e^{\alpha x} (D + a) [(D + a) u]  \\
                &= e^{\alpha x} ( D + a)^2 u 
\end{align*}
按此迭代法, 由数学归纳法推得 $p(D) = D^n$也满足上述关系.\\
而$P(D)$的一般情况为$D, D^2, D^3, ..., D^n$的无关线性组合, 
而线性组合不影响结论, 即证毕.\\
得此法则后, 下直接给出 $p(\alpha) = 0$时, 解情况:
\[ p(D) y = e^{\alpha x} \]
\[ y_p = \frac{x e^{\alpha x}}{p'(\alpha)}  \hspace*{2em} p(\alpha) = 0 \]
此时 $\alpha$是多项式$p$的单根\\
自然引出, 若$\alpha$是多项式$p$的二重根, 
即$p'(\alpha) = 0$, 解情况为?如下:
\[ y_p = \frac{x^2 e^{\alpha x}}{p''(\alpha)}  \hspace*{2em} p'(\alpha) = 0 \]
往下更高阶按此规律即可推得, 此处不再赘述.\\
下补充:\\
单根和$p(\alpha) = 0$联系:
\begin{align*}
    p(D)  &= (D - a) (D - b) \hspace*{2em} a \neq b\\
    p'(D) &= (D - b) + (D - a)  \\
    p'(a) &= a - b \neq 0 \\
    p(D) \frac{e^{\alpha x}x}{p'(\alpha)}   &= e^{\alpha x} ( D + a - b)D \frac{x}{p'(\alpha)} \\
            &= e^{\alpha x} \frac{a-b}{p'(\alpha)}\\
            &= e^{\alpha x}
\end{align*}
证毕.
